\documentclass[11pt]{article}
\input{style-pc}
\usepackage{array}

\begin{document}

\entetesujet{Sujet bac 2019 -- Physique-Chimie -- Série D}

% ============================ CHIMIE ============================
\matiere{CHIMIE}{8 points}

\partie{A : vérification des connaissances}

\noindent\textbf{1. Texte à trous.} Reproduis puis complète la phrase suivante par quatre des cinq mots ci-après : \textit{électron ; longueur ; atomes ; particules ; rayonnement}.

\textit{Lorsque les noyaux de certains $\cdots\cdots$ se désintègrent, ils émettent d'une part des $\cdots\cdots$, d'autre part un $\cdots\cdots$ électromagnétique de très courte $\cdots\cdots$ d'onde.}

\medskip
\noindent\textbf{2. Question à réponse construite.} Donne la définition d'une famille radioactive.

\medskip
\noindent\textbf{3. Questions à alternative vrai ou faux.} Réponds par vrai ou faux aux affirmations suivantes. Exemple : 3.e = vrai.
\begin{enumerate}[label=3.\alph*.]
  \item Un indicateur coloré est un acide faible ou une base faible dont la couleur de sa forme acide est différente de la couleur de sa forme basique, en solution aqueuse.
  \item Un atome d'hydrogène dans son état fondamental émet un rayonnement électromagnétique.
  \item L'élévation ébulliométrique est inversement proportionnelle à la masse molaire moléculaire du soluté.
  \item Dans les conditions normales de température et de pression, le volume molaire est $V_m = 22{,}4$ mol/L.
  \item La constante d'équilibre d'une réaction chimique équilibrée dépend des pressions partielles des réactifs et des produits.
\end{enumerate}

\partie{B : application des connaissances}
La cinétique de la réaction $\ch{2\,NO_2 \longrightarrow 2\,NO + O_2}$ a été étudiée expérimentalement à $400\,^\circ$K. On a obtenu les résultats suivants :

\begin{center}\renewcommand{\arraystretch}{1.3}
\begin{tabular}{|l|c|c|c|}
\hline
Expérience & 1 & 2 & 3 \\
\hline
$[\ch{NO_2}]$ (mol$\cdot$L$^{-1}$) & 0,85 & 1,10 & 1,6 \\
\hline
$V_0$ (mol$\cdot$L$^{-1}\cdot$s$^{-1}$) & 0,39 & 0,65 & 1,38 \\
\hline
\end{tabular}
\end{center}
\noindent où $V_0 = (\text{Vitesse})_0$.

\begin{enumerate}
  \item Déterminer l'ordre global de cette réaction ainsi que la valeur numérique de la constante de vitesse.
  \item Déduis la loi de vitesse ainsi que la loi intégrée de cette réaction.
  \item Calcule le temps de demi-réaction en considérant l'expérience 2.
  \item Au bout de combien de temps 75 \% de $[\ch{NO_2}]$ se sont-ils transformés pour l'expérience 2 ?
\end{enumerate}

% ============================ PHYSIQUE ============================
\matiere{PHYSIQUE}{12 points}

\partie{A : vérification des connaissances}

\noindent\textbf{1. Appariement.} Associe chaque élément-question de la colonne A avec un élément-réponse de la colonne B correspondante. Exemple : $A_7 = B_9$.

\begin{center}\renewcommand{\arraystretch}{1.6}
\begin{tabular}{|p{0.4\linewidth}|p{0.42\linewidth}|}
\hline
\textbf{Colonne A} & \textbf{Colonne B} \\
\hline
$A_1$ : incertitude relative & $B_1$ : nœud \\
\hline
$A_2$ : point d'amplitude maximale & $B_2$ : nombre abstrait \\
\hline
$A_3$ : point d'amplitude nulle & $B_3$ : dépend du temps \\
\hline
$A_4$ : incertitude absolue & $B_4$ : nombre concret \\
\hline
 & $B_5$ : ventre \\
\hline
\end{tabular}
\end{center}

\noindent\textbf{2. Question à réponse courte.} Donne le nom de la tension électrique qu'il faut appliquer entre l'anode et la cathode d'une cellule photoélectrique pour annuler le courant photoélectrique.

\medskip
\noindent\textbf{3. Schéma à compléter.} Dans le schéma ci-dessous, le point $M$ est un point de la périphérie d'une roue en rotation. Il effectue un mouvement circulaire uniformément accéléré.

\begin{center}
\begin{tikzpicture}[>=Stealth,scale=1]
  \draw[thick] (0,0) circle (1.9);
  \fill (0,0) circle (1.2pt) node[below right]{$O$};
  % point M
  \coordinate (M) at ({1.9*cos(130)},{1.9*sin(130)});
  \fill (M) circle (2pt) node[above left]{$M$};
  % axe rotation (Delta)
  \draw[dashed] (0.5,2.6)--(-0.5,-2.6);
  \node[right] at (0.35,1){$(\Delta)$};
  \node[below right] at (0.2,-0.3){\small axe de rotation};
  % sens du mouvement
  \draw[->,thick] (-2.5,-0.6) arc[start angle=-150,end angle=-210,radius=1.4];
  \node[left] at (-2.5,0.2){\small sens du};
  \node[left] at (-2.5,-0.1){\small mouvement};
\end{tikzpicture}
\end{center}

\noindent Recopie et complète le schéma en représentant le vecteur accélération $\vec{a}$ et le vecteur vitesse $\vec{v}$ du point $M$.

\partie{B : application des connaissances}
Un ascenseur démarre en mouvement de translation rectiligne vers le haut. Il atteint la vitesse de 6 m$\cdot$s$^{-1}$ après un parcours de 6 m. Il conserve ensuite cette vitesse sur une certaine distance $d$ puis s'arrête $2{,}5$ secondes après un parcours $d'$.

\begin{enumerate}
  \item Précise la nature du mouvement de l'ascenseur dans chacune des phases.
  \item On suspend au plafond de l'ascenseur un ressort à spires non jointives de longueur à vide $l_0 = 20$ cm et de constante de raideur $k = 30$ N$\cdot$m$^{-1}$. À l'autre extrémité de ce ressort est accroché un solide $(S)$ de masse $m = 300$ g. Au cours du mouvement de l'ascenseur, le ressort prend une longueur $L$. La masse du ressort est négligeable.
  \begin{enumerate}[label=\alph*.]
    \item Établis l'expression de la longueur $L$ du ressort en fonction de $L_0$, $m$, $k$, $g$ et de l'accélération $a$ du mouvement.
    \item Calcule la valeur numérique de la longueur $L$ au cours de chaque phase du mouvement.
  \end{enumerate}
\end{enumerate}
\noindent N.B. : on suppose qu'il ne se produit pas d'oscillations. $g = 10$ m$\cdot$s$^{-2}$.

\partie{C : résolution d'un problème}
Le but de ce problème est de déterminer la période $T$ d'un pendule conique. On dispose d'un ressort vertical à spires non jointives de longueur à vide $L_0 = 20$ cm. On accroche un solide $S$ de masse $m = 200$ g à l'extrémité inférieure du ressort. À l'équilibre, la longueur du ressort est $L_1 = 30$ cm.

\begin{minipage}{0.6\linewidth}
\begin{enumerate}
  \item Détermine la constante de raideur $K$ du ressort.
  \item Ce ressort est fixé par son extrémité supérieure en un point $O$ d'un axe vertical $(\Delta)$. L'ensemble est mis en rotation autour de l'axe $(\Delta)$ grâce à un moteur. Le solide $(S)$ décrit alors un cercle dans un plan horizontal et la direction du ressort fait un angle $\alpha = 30^\circ$ avec l'axe $(\Delta)$.
  \begin{enumerate}[label=\alph*.]
    \item Représente toutes les forces qui s'appliquent sur le solide $(S)$.
    \item Calcule :
    \begin{enumerate}[label=b.\arabic*.]
      \item La longueur $L_2$ du ressort lors de ce mouvement.
      \item La vitesse angulaire $\omega$ de rotation de l'ensemble.
    \end{enumerate}
  \end{enumerate}
  \item Calcule enfin la période $T$ de ce pendule conique.
\end{enumerate}
\end{minipage}\hfill
\begin{minipage}{0.36\linewidth}
\centering
\begin{tikzpicture}[>=Stealth,scale=1]
  \coordinate (O) at (0,2.7);
  \coordinate (S) at (2.3,-0.3);
  % axe Delta
  \draw[dashed] (0,3)--(0,-1.4);
  \node[right] at (0,-1.2){$(\Delta)$};
  \fill (O) circle (1pt) node[above]{$O$};
  % ressort (coil) de O vers S
  \draw[decorate,decoration={coil,aspect=0.6,segment length=5pt,amplitude=4pt}] (O)--(S);
  % solide
  \draw[fill=gray!30,rotate around={-50:(S)}] ($(S)+(-0.18,-0.12)$) rectangle ($(S)+(0.18,0.12)$);
  \node[right] at (S){$(S)$};
  % trajectoire (cercle horizontal en ellipse)
  \draw (0,-0.3) ellipse (2.3 and 0.55);
  % angle alpha
  \draw (0,2.1) arc[start angle=270,end angle={270+38},radius=0.6];
  \node at (0.32,2.1){$\alpha$};
\end{tikzpicture}
\end{minipage}
\noindent\\[2pt] On donne : $g = 10$ m/s$^2$.

\end{document}
